% =====================================================================
% Bases de dados e grafos sintéticos --- Brunno (Frente 1)
% Só o Brunno edita este arquivo.
% =====================================================================
\section{Bases de dados}\label{sec:bases}

Usamos quatro redes públicas e anonimizadas. As funções de \codigo{leitores/bases.py} leem os
arquivos e descartam laços (arestas de um nó para ele mesmo). As quatro são não direcionadas e conexas, e a
Tabela~\ref{tab:bases} as resume com números calculados pela ferramenta; os testes dos leitores
conferem que os tamanhos lidos são os publicados. Nas três maiores, o grau é muito desigual: a
mediana fica bem abaixo da média, e o maior grau passa de 20 vezes a média na Twitch PT e no
ego-Facebook, uma cauda longa comum em redes sociais~\cite{newman2018}.

\begin{table}[htbp]
\centering
\caption{Bases de dados: $n$ nós, $m$ arestas, densidade $2m/(n(n-1))$ e graus médio, mediano e
máximo, calculados pela ferramenta.}
\label{tab:bases}
\small
\setlength{\tabcolsep}{4.5pt}
\begin{tabular}{@{}l r r r r r r l@{}}
\toprule
Base & $n$ & $m$ & Densidade & \multicolumn{3}{c}{Grau} & Uso \\
\cmidrule(lr){5-7}
 & & & & médio & mediano & máximo & \\
\midrule
Twitch PT & 1.912 & 31.299 & 0,0171 & 32,7 & 17 & 767 & Influenciadores \\
ego-Facebook & 4.039 & 88.234 & 0,0108 & 43,7 & 25 & 1.045 & Desempenho \\
Network of Thrones & 303 & 1.008 & 0,0220 & 6,7 & 4 & 59 & Demonstração \\
Karate Club & 34 & 78 & 0,1390 & 4,6 & 3 & 17 & Corretude \\
\bottomrule
\end{tabular}
\end{table}

A \textbf{Twitch PT} é a rede de amizades mútuas entre streamers que transmitem em português,
coletada por \cite{rozemberczki2021} e distribuída pelo SNAP~\cite{leskovec2014snap}. Cada nó traz
atributos reais: as visualizações do canal, que servem de validação externa para o ranking de
influenciadores, se ele é parceiro da Twitch (279 dos 1.912), se tem conteúdo adulto e há quantos
dias existe. O \textbf{ego-Facebook} une dez redes ego de amizades~\cite{mcauley2012} e, por ser a maior,
serve para medir o desempenho. Na \textbf{Network of
Thrones}, dois personagens do terceiro livro de \emph{As Crônicas de Gelo e Fogo} se ligam quando os
nomes aparecem a até 15 palavras um do outro, e o peso, de 3 a 95, conta quantas vezes isso
acontece~\cite{beveridge2016}; ela serve à demonstração, porque o público reconhece os personagens.
Por fim, o \textbf{Karate Club} é o clube que se dividiu em dois grupos de 17
membros~\cite{zachary1977}, e os seus valores conhecidos servem para provar a corretude.

\paragraph{Grafos sintéticos.} Para testar os algoritmos em tamanhos controlados e medir o
desempenho, \codigo{geradores/modelos.py} gera três modelos clássicos, com uma \codigo{semente} que
repete o grafo. No de \textbf{Erdős--Rényi}~\cite{erdos1959}, cada par de nós se liga com
probabilidade $p$, e o gerador pula direto para o próximo par sorteado, em $O(n + m)$, em vez de
sortear os $n(n-1)/2$ pares~\cite{batagelj2005}. No de \textbf{Barabási--Albert}~\cite{barabasi1999},
cada nó novo se liga a $m$ nós já existentes, com chance proporcional ao grau, o que cria nós muito
conectados, como nas redes reais. No de \textbf{Watts--Strogatz}~\cite{watts1998}, cada nó de um anel
se liga aos $k$ mais próximos e cada aresta é religada com probabilidade $p$, o que mantém o
agrupamento alto e encurta as distâncias. A interface usa os mesmos geradores.
